\songtitle{\textjapanese{夏の終わりだった} (Natsu no owari datta)}

\tag*{2026}\tag{electroswing}\tag{japanese-lyrics}\tag{final-del-verano}
\begin{notes}
Shorter electroswing adaptation into Japanese of \emph{El final del verano} (page \pageref{sec:verano}), by Claude; compare also \emph{\japanesefont 湖畔フェス} for the full-length Japanese version.
\end{notes}
\begin{style}
Electroswing with a high-tempo dance beat at 124 BPM in F minor, The track features a prominent brass section with trumpets and trombones playing syncopated staccato stabs, A synthesized upright bass provides a walking bassline that interacts with a four-on-the-floor kick drum and a crisp, high-frequency snare, A bright, processed female vocal performs repetitive melodic hooks with heavy compression and a slight slapback delay, Glitchy record scratches and vinyl crackle effects are layered over the percussion, A swing-style piano plays rhythmic chords on the off-beats, The arrangement uses sudden filter sweeps and momentary silences for rhythmic emphasis
\end{style}
\begin{lyrics}
\songpart{Intro}
\annotation{four-on-the-floor kick drum, vinyl crackle, staccato brass stabs}\\
\begin{japanese}
夏の終わりだった\\
\end{japanese}
\annotation{record scratch}\\
\begin{japanese}
君を見た瞬間\\
\end{japanese}
\annotation{full band enters, walking synth bass, swing piano}\\
\begin{japanese}
湖畔のフェスだった\\
理性なんて失くした瞬間
\end{japanese}

\songpart{Chorus}
\annotation{processed female vocals, bright brass section}\\
\begin{japanese}
夏の終わりだった\\
君に落ちた瞬間
\end{japanese}

\songpart{Outro}
\annotation{filter sweep on brass}\\
\begin{japanese}
そして君は\\
一瞬で 私の全てになった\\
君は私を
\end{japanese}
\end{lyrics}

